\begin{figure}[htbp]
  \centering
  \begin{tikzpicture}[x=7mm,y=-7mm,alt={A five by three output is partitioned into tiles of edge three. The top tile has three rows and three columns, and the lower clipped tile has two rows and three columns. Along a separate seven-term reduction axis, blocks cover indices zero through two, three through five, and six alone. The last output tile and last reduction block are clipped independently.}]
    \fill[black!14] (0,3) rectangle (3,5);
    \foreach \x in {0,...,3} {\draw[black!40] (\x,0)--(\x,5);}
    \foreach \y in {0,...,5} {\draw[black!40] (0,\y)--(3,\y);}
    \draw[thick] (0,0) rectangle (3,5) (0,3)--(3,3);
    \node[font=\small,anchor=south] at (1.5,-.15) {$M=5,\ N=3$};
    \node[font=\small,anchor=north] at (1.5,5.15) {Clipped tile: $2\times3$};
    \foreach \x in {0,...,6} {\node[font=\small] at (6+\x+.5,1.5) {\x};}
    \foreach \x in {0,...,7} {\draw[black!40] (6+\x,1)--(6+\x,2);}
    \draw[thick] (6,1) rectangle (13,2) (9,1)--(9,2) (12,1)--(12,2);
    \node[font=\small,anchor=south] at (9.5,.7) {$K=7,\ T=3$};
    \node[font=\small,anchor=north] at (9.5,2.3) {Reduction lengths: 3, 3, 1};
    \node[font=\small,align=left,anchor=west] at (6,4) {Every output receives all seven terms;\\the final reduction block adds one term.};
  \end{tikzpicture}
  \caption{Output clipping and reduction clipping are separate bounds. The final output tile can contain several cells while its final reduction interval has only one term.}
  \label{fig:revision-measurement-tiles}
\end{figure}
